% Generated draft: requires mathematical and source review.
\documentclass[11pt]{article}
\usepackage[margin=25mm]{geometry}
\usepackage{fontspec}
\setmainfont{Latin Modern Roman}
\usepackage{amsmath,amssymb,mathtools,unicode-math}
\setmathfont{Latin Modern Math}
\usepackage{enumitem,needspace,xurl,hyperref}
\setlength{\parindent}{0pt}
\setlength{\parskip}{5pt}
\setlength{\emergencystretch}{3em}
\newcommand{\N}{\mathbb{N}}
\newcommand{\Z}{\mathbb{Z}}
\newcommand{\Q}{\mathbb{Q}}
\newcommand{\R}{\mathbb{R}}
\newcommand{\C}{\mathbb{C}}
\newcommand{\F}{\mathbb{F}}
\newcommand{\A}{\mathbb{A}}
\newcommand{\PP}{\mathbb{P}}
\newcommand{\E}{\mathbb{E}}
\newcommand{\eps}{\varepsilon}
\newcommand{\dd}{\,\mathrm{d}}
\begin{document}
\section*{UnsolvedMath ID: 20003308}
\textbf{Source problem code:} catalogue entry
\subsection{Definitions and mathematical statement}
The assigned extract is transitional prose: it says that Conjecture 29 is a special case of a general Cartan-invariant closed-set conjecture [130], then introduces another conjecture of G. Margulis without stating it. Companion records identify the two different assertions. For
\[
 X_3=\mathrm{SL}_3(\R)/\mathrm{SL}_3(\Z),\qquad
 A=\{\operatorname{diag}(a_1,a_2,a_3):a_i\ne0,\ a_1a_2a_3=1\},
\]
Conjecture 29 says that every relatively compact $A$-orbit is closed. Conjecture 30 says that for every compact $K\subset X_3$, only finitely many closed $A$-orbits are contained in $K$. The latter concerns containment of whole orbits, not merely their meeting $K$. This context record does not prove either assertion or an implication between them.

The attached aid reports logical independence in abstract dynamics. For a continuous action of $\R^2$ on a metric space, consider properties
\[
 \text{every relatively compact orbit is closed},\qquad
 \text{each compact set contains only finitely many closed orbits}.
\]
Two faithful actions on compact metric spaces separate them. First, on $[0,1]^2$, let
\[
 \psi_s(x)=\frac{e^s x}{1-x+e^s x},\qquad
 (s,t)\cdot(x,y)=(\psi_s(x),\psi_t(y)).
\]
Exactly the four corner orbits are closed, so compact-set finiteness holds, while the interior orbit is relatively compact and not closed. Second, on $\mathbb T^2\times[1,2]$, let
\[
 (s,t)\cdot(x,y,r)=(x+rs,y+rt,r),
 \qquad(x,y)\in\mathbb T^2.
\]
Every orbit is a full closed torus fiber, but the whole compact space contains infinitely many such orbits. The varying scale $r$ makes this $\R^2$-action faithful. These examples show that neither abstract property alone implies the other; they are not counterexamples in $X_3$.

The aid additionally reports a root-unipotent escape lemma and a conditional orbit-closure amplification proposition isolating the extra homogeneous input behind an An--Weiss compact-minimal-set reduction. Its logical structure requires an amplification hypothesis forcing a suitable nonclosed compact invariant configuration to contain a complete root-unipotent orbit, together with an escape statement excluding that containment. No unconditional amplification or implication between the arithmetic conjectures is supplied, and full hypotheses and references for this auxiliary reduction are absent from the summary.
\subsection{Short English statement}
This transition links two distinct Margulis conjectures: closing bounded diagonal orbits and finiteness of closed orbits in compact sets. The supplied compact-space examples show those properties are independent in abstract dynamics; extra homogeneous hypotheses are needed for any arithmetic implication.
\subsection{Sources}
\begin{itemize}
\item American Institute of Mathematics, \emph{Emerging applications of measure rigidity}, transition between Conjectures 29 and 30, attributed to G. Margulis and referring to [130], \url{https://aimath.org/WWN/measrigid/measrigid.pdf}. The full bibliography entry for [130] is not supplied.
\item Companion local records in \texttt{aim-other-notes.json}: permanent ID $20003307$, source code \texttt{AIM-OTHER-0114}, \texttt{notes[113]}, supplies Conjecture 29; permanent ID $20003309$, \texttt{AIM-OTHER-0116}, \texttt{notes[115]}, supplies Conjecture 30.
\item Local UnsolvedMath registry, permanent ID $20003308$, source code \texttt{catalogue entry}, canonical \texttt{aim-other-notes.json}, \texttt{notes[114]}, dated 14 August 2026. Research attempt 1: context only, ``not a problem'' at run, proved rigor, low novelty confidence; machine-generated and not peer reviewed. The summary names An--Weiss but gives no complete citation.
\end{itemize}
\end{document}
